\documentclass[a4paper,12pt]{article}
\usepackage[utf8]{inputenc}
\usepackage[spanish]{babel}
\usepackage{amsmath}
\usepackage{siunitx}
\usepackage{geometry}
\usepackage{booktabs}
\usepackage{graphicx}
\usepackage{caption}
\usepackage{circuitikz}

\begin{document}

% Carátula
\begin{titlepage}
    \centering
    \vspace*{2cm}
    \textbf{\Large Instituto Tecnológico de Buenos Aires} \\
    \vspace{0.5cm}
    \textbf{\large Física III} \\
    \vspace{1cm}
    \textbf{\LARGE Trabajo Práctico N°3: Circuitos de Corriente Alterna} \\
    \vspace{1.5cm}
    \begin{tabular}{ll}
        \textbf{Grupo:} & 15 \\
        \textbf{Integrantes:} & Nicolás Valentín Arias (Legajo 62272) \\
                              & Julián Yaregui Murphy (Legajo 64616) \\
                              & Guillermo Arredondo (Legajo 58082) \\
                              & Octavio Zacagnino (Legajo 64255) \\
        \textbf{Fecha de realización:} & 11/06/2025 \\
        \textbf{Fecha de entrega:} & 18/06/2025 \\
    \end{tabular}
    \vspace{2cm}
    \begin{tabular}{|p{12cm}|}
        \hline
        \textbf{Observaciones del docente:} \\
        \vspace{2cm} \\
        \hline
    \end{tabular}
    \vspace{1cm}
    \begin{tabular}{p{6cm}}
        \textbf{Firma del docente:} \\
        \vspace{1cm} \\
        \hrule
    \end{tabular}
\end{titlepage}

% Contenido del informe
\section{Objetivos y Resumen}
% Estableciendo los objetivos del experimento
Este trabajo práctico tuvo como objetivos principales:
\begin{itemize}
    \item Analizar el comportamiento de un circuito RLC serie al variar la frecuencia de la fuente, determinando la frecuencia de resonancia \( f_0 \), la inductancia \( L \) y el factor de calidad \( Q \).
    \item Estudiar circuitos rectificadores de media onda y onda completa, evaluando el efecto del capacitor de filtrado en la señal de salida y comparando su eficiencia energética.
\end{itemize}

% Resumiendo los resultados obtenidos
En el experimento, se midió la corriente en un circuito RLC serie en función de la frecuencia, identificando una frecuencia de resonancia aproximada de \SI{4000}{\hertz}, con una inductancia calculada de \SI{11}{\milli\henry} y un factor de calidad \( Q \approx 0.645 \), indicando alta amortiguación. En la sección de rectificadores, se observó que el capacitor de filtrado reduce el rizado, y el rectificador de onda completa mostró mayor eficiencia. Estos resultados tienen aplicaciones en filtros electrónicos y fuentes de alimentación.

% Agregando la imagen del osciloscopio
\begin{figure}[h]
    \centering
    \includegraphics[width=0.8\textwidth]{imagen_osciloscopio.jpeg}
    \caption{Oscilograma mostrando la caída de tensión en la resistencia a \SI{5.6}{\volt}, utilizada para determinar la frecuencia de resonancia.}
\end{figure}

\section{Introducción Teórica}
% Introduciendo los conceptos teóricos fundamentales
\subsection{Circuitos RLC Serie}
Un circuito RLC serie consta de una resistencia \( R \), un inductor \( L \) y un capacitor \( C \) conectados en serie a una fuente de tensión alterna \( V(t) = V_0 \sin(\omega t) \). La impedancia total del circuito se expresa como:
\[
Z = R + j \left( \omega L - \frac{1}{\omega C} \right)
\]
donde \( \omega = 2\pi f \) es la frecuencia angular. La corriente alcanza su máximo cuando la reactancia inductiva (\( \omega L \)) iguala la reactancia capacitiva (\( \frac{1}{\omega C} \)), lo que ocurre en la frecuencia de resonancia:
\[
\omega_0 = \frac{1}{\sqrt{L C}} \quad \Rightarrow \quad f_0 = \frac{1}{2\pi \sqrt{L C}}
\]
En resonancia, \( Z = R \), y la corriente máxima es \( I_0 = V_0 / R \). El factor de calidad \( Q \), que mide la selectividad del circuito, se define como:
\[
Q = \frac{f_0}{\Delta f} = \frac{1}{R} \sqrt{\frac{L}{C}}
\]
donde \( \Delta f = f_2 - f_1 \), siendo \( f_1 \) y \( f_2 \) las frecuencias donde la corriente cae a \( I_0 / \sqrt{2} \).

% Agregando la imagen del circuito
\begin{figure}[h]
    \centering
    \includegraphics[width=0.8\textwidth]{imagen_circuito.jpeg}
    \caption{Configuración física del circuito RLC serie utilizado en el experimento.}
\end{figure}

% Explicando el circuito RLC paralelo y el circuito trampa
\subsection{Circuitos RLC Paralelo y Circuitos Trampa}
En un circuito RLC paralelo, los componentes están conectados en paralelo, y la admitancia total es:
\[
Y = \frac{1}{R} + j \left( \omega C - \frac{1}{\omega L} \right)
\]
En resonancia, la admitancia es mínima, bloqueando ciertas frecuencias. Esto se conoce como "circuito trampa", usado para eliminar frecuencias no deseadas, como interferencias en receptores de radio ([Circuitos RLC](https://es.wikipedia.org/wiki/Circuito_RLC)).

% Detallando los rectificadores
\subsection{Circuitos Rectificadores}
Los rectificadores convierten corriente alterna (AC) en corriente continua (DC). En un rectificador de media onda, un diodo permite el paso de un semiciclo, produciendo una salida pulsante. En un rectificador de onda completa, un puente de diodos utiliza ambos semiciclos, duplicando la frecuencia de la salida. Un capacitor en paralelo con la carga reduce el rizado, con una tensión de rizado aproximada:
\[
V_{\text{ripple}} \approx \frac{I_{\text{carga}}}{f C}
\]
donde \( f \) es la frecuencia de la fuente y \( C \) es la capacitancia.

\section{Instrumentos y Elementos Empleados}
% Listando los instrumentos para el circuito RLC
\subsection{Circuito RLC Serie}
\begin{itemize}
    \item Resistencia \( R \) de valor conocido.
    \item Capacitor \( C = \SI{144e-9}{\farad} \).
    \item Inductor \( L \) (valor a determinar).
    \item Generador de funciones (señal sinusoidal, amplitud constante).
    \item Multímetro para medir corriente (a través de la tensión en \( R \)).
\end{itemize}

\subsubsection{Procedimiento 2.1}
En esta primera etapa del experimento, se dispuso el circuito RLC serie y se varió la frecuencia de la fuente, midiendo la caída de potencial sobre el resistor con el osciloscopio para identificar la frecuencia de resonancia.

\subsection{Circuitos Rectificadores}
\begin{itemize}
    \item Fuente de alterna (\SI{12}{\volt} RMS, \SI{50}{\hertz}).
    \item Diodos para rectificación.
    \item Resistencias de carga (\SI{100}{\ohm} y un cooler).
    \item Capacitores de filtrado (valores a especificar).
    \item Osciloscopio y multímetro para medir tensiones.
\end{itemize}

% Incluyendo un esquema del circuito
\begin{figure}[h]
    \centering
    \begin{circuitikz}
        \draw
        (0,0) to[sinusoidal voltage source] (0,2)
        to[R=$R$] (2,2)
        to[L=$L$] (4,2)
        to[C=$C$] (6,2)
        to[short] (6,0)
        to[short] (0,0);
    \end{circuitikz}
    \caption{Esquema de un circuito RLC serie.}
\end{figure}

\section{Datos Obtenidos: Discusión y Cálculo de Errores}
% Presentando los datos del circuito RLC
\subsection{Circuito RLC Serie}
Se midió la corriente en función de la frecuencia, obteniendo los siguientes datos:

\begin{table}[h]
    \centering
    \begin{tabular}{cc}
        \toprule
        Frecuencia (\si{\hertz}) & Corriente (\si{\ampere}) \\
        \midrule
        3040 & 0.0456416 \\
        4365 & 0.0167164 \\
        5230 & 0.0561194 \\
        6040 & 0.0501442 \\
        8072 & 0.0385074 \\
        9325 & 0.0331340 \\
        13400 & 0.0244000 \\
        \bottomrule
    \end{tabular}
    \caption{Mediciones de corriente vs. frecuencia.}
\end{table}

% Analizando la frecuencia de resonancia
La frecuencia de resonancia se estimó en \( f_0 = \SI{4000}{\hertz} \), posiblemente por interpolación, aunque el máximo de corriente aparece en \SI{5230}{\hertz} (\SI{0.0561}{\ampere}). Usando \( f_0 = \SI{4000}{\hertz} \) y \( C = \SI{144e-9}{\farad} \):
\[
L = \frac{1}{(2\pi f_0)^2 C} = \frac{1}{(2\pi \cdot 4000)^2 \cdot 144 \cdot 10^{-9}} \approx \SI{0.011}{\henry} = \SI{11}{\milli\henry}
\]
Este valor corrige el reportado previamente (\SI{8.16}{\milli\henry} o \SI{3.16}{\milli\henry}).

% Calculando el factor de calidad
El factor de calidad se calculó como:
\[
Q = \frac{f_0}{f_1 - f_2}
\]
Con \( f_1 = \SI{7100}{\hertz} \) y \( f_2 = \SI{900}{\hertz} \):
\[
Q = \frac{4000}{7100 - 900} \approx 0.645
\]
Este \( Q \) bajo sugiere alta amortiguación, posiblemente por una resistencia elevada. Teóricamente:
\[
Q = \frac{1}{R} \sqrt{\frac{L}{C}}
\]
Sin \( R \), no se calcula numéricamente, pero el valor medido indica un circuito poco selectivo.

% Discutiendo posibles errores
Posibles fuentes de error incluyen:
\begin{itemize}
    \item Imprecisiones en la medición de frecuencia o corriente.
    \item Resistencia parásita del inductor.
    \item Variaciones en \( C \) debido a tolerancias del componente.
\end{itemize}

% Presentando los datos de los rectificadores
\subsection{Circuitos Rectificadores}
\textbf{Media Onda:} Se midió la tensión de salida:
- Sin capacitor: \( V_{\text{out}} \) pulsante, valores a especificar.
- Con capacitor: \( V_{\text{out}} \) más constante, rizado reducido (mediciones pendientes).

\textbf{Onda Completa:} Similar, con mayor \( V_{\text{out}} \) promedio y menor rizado.

% Sugiriendo mediciones específicas
Se recomienda medir:
- Tensión pico y RMS sin capacitor.
- Tensión promedio y rizado con diferentes valores de \( C \).
- Comparar con \( V_{\text{ripple}} \approx \frac{I_{\text{carga}}}{f C} \).

\section{Análisis de los Resultados Obtenidos}
% Analizando el circuito RLC
\subsection{Circuito RLC Serie}
La frecuencia de resonancia estimada (\SI{4000}{\hertz}) y \( L = \SI{11}{\milli\henry} \) son consistentes con la teoría, aunque el máximo de corriente en \SI{5230}{\hertz} sugiere una posible discrepancia en \( f_0 \). El \( Q = 0.645 \) indica un circuito con alta amortiguación, lo que reduce su selectividad, útil en aplicaciones donde se desea un ancho de banda amplio.

% Analizando los rectificadores
\subsection{Circuitos Rectificadores}
El capacitor de filtrado reduce el rizado al almacenar carga durante los picos y liberarla durante los valles. El rectificador de onda completa mostró mayor eficiencia debido a la utilización de ambos semiciclos. Respuestas a preguntas:
\begin{enumerate}
    \item Mayor capacitancia reduce el rizado, según \( V_{\text{ripple}} \approx \frac{I_{\text{carga}}}{f C} \).
    \item Se llama "capacitor de filtrado" porque elimina componentes alternas.
    \item La velocidad del cooler aumenta con mayor \( V_{\text{out}} \) constante.
    \item El multímetro mostró mayor \( V_{\text{out}} \) con capacitor.
    \item El rizado es mayor en media onda sin capacitor.
\end{enumerate}

\section{Aplicaciones Prácticas}
% Describiendo aplicaciones de los circuitos RLC
Los circuitos RLC serie son fundamentales en filtros pasa-banda, usados en sintonizadores de radio para seleccionar frecuencias específicas ([Circuitos RLC](https://www.electricity-magnetism.org/es/circuitos-rlc-propiedades-aplicaciones-y-ejemplos/)). Los circuitos RLC paralelo, o "circuitos trampa", eliminan frecuencias no deseadas, como interferencias en receptores de radio ([Circuito Trampa](https://es.frwiki.wiki/wiki/Circuit_RLC)). En la industria médica, se usan en equipos para filtrar señales ([Aplicaciones RLC](https://prezi.com/tngwx9qh_yqq/aplicaciones-de-circuitos-rlc/)).

% Describiendo aplicaciones de los rectificadores
Los rectificadores son esenciales en fuentes de alimentación para dispositivos electrónicos, como cargadores y adaptadores, convirtiendo AC a DC con alta eficiencia ([Rectificadores](https://www.ingenierizando.com/electronica/rectificador/)).

\section{Conclusión}
% Resumiendo los hallazgos y su relevancia
Este trabajo práctico permitió estudiar el comportamiento de circuitos RLC serie y rectificadores. La frecuencia de resonancia (\SI{4000}{\hertz}) y la inductancia (\SI{11}{\milli\henry}) confirman la teoría, aunque el bajo \( Q = 0.645 \) sugiere alta amortiguación, posiblemente por resistencia parásita. Los rectificadores demostraron que el capacitor de filtrado reduce el rizado, y el de onda completa es más eficiente. Estos conceptos son fundamentales en aplicaciones como sintonizadores de radio, filtros de interferencia y fuentes de alimentación. Para futuros experimentos, se recomienda medir más puntos alrededor de \( f_0 \) y especificar datos de los rectificadores para un análisis más completo.

\end{document}